\chapter{What Goes Into a Logic}
\label{sec:oslf-inputs}
% ===================================================================

\begin{remark}[The state of this chapter]
\label{rmk:oslf-status}
This chapter is where the constructions that feed the OSLF
functor --- Operational Semantics in Logical Form, the machinery
that manufactures a modal logic out of a rewrite theory rather
than bolting one on beside it --- are gathered.
At present it contains one of them, in full: the construction of
\emph{space} from term structure, which the rest of this book uses
more than any other and which has until now been assumed rather
than built.
The remaining ingredients --- the SOS presentation the functor
consumes, the adjunction it factors through, and the sense in
which its output is the internal language of the theory it came
from --- belong here too, and are not yet written.
The reader should treat the chapter as load-bearing but
unfinished, and should not conclude from its brevity that the
OSLF construction is brief.
\end{remark}

\section{The agent needs a \emph{where}}

Everything so far has been a theory of state and of evolution.
A grammar says which configurations are possible, an equational
quotient says which of them are the same, and a collection of
rewrites says how one becomes another.
That is enough to have a time: Chapter~\ref{sec:causal} read time
off paths through the synchronization tree, and
Chapter~\ref{sec:reversibility} read reversible time off paths
that carry their own history.

It is not yet enough to have a \emph{space}.
An agent working out its situation needs to know not only what is
happening but \emph{where} --- where the event is, where its
environment is relative to itself, whether two things that just
happened happened near each other or far apart.
And the agent cannot be handed a space, because there is nothing
to hand it one.
It has its theory of state and evolution and nothing else.
Whatever space it has, it must read off what it already has.

The good news is that what it already has is not featureless.
Terms have structure.
They are trees, and once the equational quotient is imposed,
graphs: constructors at the nodes, subterms in the argument
positions.
That structure is the only candidate for a notion of space
available to an agent confined to its own theory, and it is also,
as it happens, a good one.
It is combinatorial, so it needs no external manifold or metric to
make sense of.
It carries adjacency for free: two argument positions of the same
constructor are next to each other.
It supports locality: a rewrite at one place leaves the
surrounding structure alone.
And it is the structure the agent's hypotheses already speak
about, so nothing has to be added to the hypothesis language to
let the agent say where.

\section{One-holed contexts, and the derivative of a type}

Chapter~\ref{sec:lts-hml} already introduced contexts: terms with
a distinguished hole, written $K[-]$, with $K[P]$ the result of
plugging $P$ into it.
There the interest was in \emph{minimal} contexts, because minimal
contexts are what make a bisimulation a congruence.
Here the interest is in contexts as a family --- in the type of
all of them.

That type has a striking description, due in its modern form to
Huet and to McBride~\cite{Huet97,McBride08}.
Huet observed that navigating a tree while retaining the ability
to rebuild it requires carrying the surroundings along as a
first-class object, and called the resulting data structure a
\emph{zipper}.
McBride identified what the zipper is: the type of one-holed
contexts over a type $T$ is the formal \emph{derivative} $\partial
T$, computed by exactly the rules one learns in a first calculus
course, read type-theoretically.
The derivative of a sum is the sum of the derivatives; the
derivative of a product obeys the Leibniz rule; the derivative of
a recursive type unfolds by the chain rule.

\begin{definition}[The context type]
\label{def:dT}
Let $\terms(\GSLT)$ be the term type of a GSLT $\GSLT$, presented
as a polynomial functor over its signature.
Its \emph{context type} is the formal derivative
$\partial \terms(\GSLT)$, whose inhabitants are the one-holed
contexts $K[-]$ of Chapter~\ref{sec:lts-hml}.
\end{definition}

The derivative is not an analogy here.
It is the same calculation, and it is worth pausing on the fact
that the rule which tells you the derivative of $x^n$ is $nx^{n-1}$
is, read as a statement about types, the observation that an
$n$-tuple has exactly $n$ places you can put a hole in and an
$(n-1)$-tuple left over when you do.
This is developed at length elsewhere~\cite{spacecalc}; what is
needed here is only the type and the plugging operation.

\section{A context is a shape; it is not a place}

It is tempting to stop here and call a context a location.
The temptation should be resisted, and the reason is instructive.

Consider the context $\lambda x.[-]$.
Does it locate the subterm $N$ inside $\lambda x.(M\,N)$?
It does not, or rather it does not do so uniquely: the same
context sits equally well around the subterm $P$ inside
$\lambda x.(M\,(N\,P))$, and around indefinitely many others.
A context tells you what the surroundings look like.
It does not tell you which surroundings, of which term, you are in
the middle of.
It is a shape, and a shape is not a place.

What locates is the \emph{pair}.

\begin{definition}[Splitting]
\label{def:splitting}
A \emph{splitting} of a term $s \in \terms(\GSLT)$ is a pair
$(K, t) \in \partial\terms(\GSLT) \times \terms(\GSLT)$ with
$K[t] = s$.
Write $\Split(\GSLT) = \partial\terms(\GSLT) \times \terms(\GSLT)$ for the
type of all splittings.
The \emph{plug map}
\[
  \Plug : \Split(\GSLT) \longrightarrow \terms(\GSLT),
  \qquad (K,t) \mapsto K[t]
\]
sends a splitting to the term it reconstructs, forgetting where
the cut was.
A \emph{location in $s$} is a point of the fibre $\Plug^{-1}(s)$.
\end{definition}

So a location is not a thing found inside a term.
It is a way of cutting the term in two, and the location \emph{is}
the cut.

\begin{remark}[Dedekind's move]
\label{rmk:dedekind}
The pattern is one mathematics has used before, and it is worth
naming because it makes the definition feel less like a
technicality.
Dedekind does not locate $\sqrt{2}$ by exhibiting a rational.
He locates it by partitioning $\mathbb{Q}$ into two halves, with
$\sqrt{2}$ the gap between them.
Conway's surreal numbers make the move constitutive: every number
\emph{is} a pair $(L \mid R)$ of left and right sets, and is
nothing but the gap they specify~\cite{ONAG}.
In each case neither half names the point; the partition does.
The analogy is structural rather than identifying --- a Dedekind
cut splits a densely ordered set and a splitting splits a
tree-shaped one, with a plug relation between the halves that the
ordered case has no use for --- but the shape of the idea is the
same.
A place is a way of dividing a whole, and the division is the
place.
\end{remark}

\section{Fire: the events available now}

Not every cut is interesting.
Most of them are inert: nothing can happen there, because nothing
at that cut matches the left-hand side of any rule.
The interactive structure of Chapter~\ref{sec:gslt} tells us which
cuts are not inert.

Recall that in an interactive GSLT the base rewrites --- those
with no hypotheses above the line, the ground events of the
theory --- all share a distinguished family $K$ of head
constructors.
In the lambda calculus that family is application; in the $\pi$-
and rho calculi it is parallel composition.
A splitting is where something can happen exactly when its
subterm side exposes such a redex.

\begin{definition}[The bundle of available events]
\label{def:fire}
$\Fire(\GSLT) \subseteq \Split(\GSLT)$ is the sub-bundle of
splittings $(K,t)$ for which $t$ matches the left-hand side of
some base rewrite of $\GSLT$, with $K$ the context in which that
rewrite is to fire.
A point of $\Fire(\GSLT)$ is an \emph{available event}: an
interaction that can happen at this cut, in this state, now.
\end{definition}

$\Fire$ is much smaller than $\Split$.
Events are sparse; at any given moment almost every way of cutting
a state in two is a way of cutting it at a place where nothing is
going on.
Sparseness is not a defect of the construction but the content of
it: it is what makes the difference between a state in which
anything might happen and a state in which specific things might.

Two consequences follow immediately, and they are the reason to
have bothered.

First, the fibre of $\Fire$ over a state $s$ --- the set of
available events at $s$ --- is the discrete analog of the
tangent space at $s$.
It is the set of directions in which the state can move.
A path through the synchronization tree, in which each step is the
firing of a base rule at some available cut, is then the discrete
analog of an integral curve of a vector field.
Time, which Chapter~\ref{sec:causal} obtained from paths, is
motion through this structure.

Second, the program/environment cut is no longer a metaphor.
In an interactive GSLT with symmetric head constructors, we said
that which side is the program and which the environment is a
perspectival fact rather than a fact about the situation.
We can now say what the perspective is: it is a splitting.
The cut tells you where the program --- the subterm the base rule
consumes --- meets the environment --- the context in which the
consuming happens.
Location and interaction are the same notion viewed at different
resolutions.

\section{Space and time from one object}

Put the two together.

\begin{quote}
Time is the path through the discrete tangent structure.
Space is the structure of cuts the path passes through.
\end{quote}

Both are read off the same object: the term structure of states,
differentiated and then restricted to the events that are
actually available.
Neither is prior to the other and neither is imported.
They share an origin, and the origin is the compression of causal
structure into rule form that a GSLT already is.

This is a substantial departure from how physical theories are
normally set up, and the departure is worth stating plainly
because the rest of this book depends on it.
In the standard picture, space is given --- a manifold, a lattice,
a Hilbert space --- and the dynamics is then defined \emph{on}
that pre-given arena.
Here there is no arena.
The derivative of the term type produces one out of the theory of
computation, and the events restrict it to the part of it that
matters.
The dynamics does not live in a space; it makes one.

\section{What is still missing}

Two things, and both are wanted before the OSLF construction can
be given in full.

The first is a metric.
Everything above is combinatorial.
The cuts are there, adjacency is there, locality is there, but
nothing yet says that two states which behave almost identically
are \emph{near} each other, and that is the notion of nearness a
predicting agent actually needs.
The stratified distance of Chapter~\ref{ch:sci} is one answer,
built for a scientist with a budget; whether it is the answer that
falls out of this construction, or merely one compatible with it,
is open.

The second is the functor itself.
What has been built here is the raw material: a type of locations,
a sub-bundle of events, and a plug map relating them to states.
OSLF consumes material of this kind and returns a modal logic
whose formulae characterize behavior up to bisimulation.
Chapter~\ref{sec:lts-hml} gave the logic; it did not give the
manufacture.
Until the manufacture is written down, the reader is entitled to
regard the logic as postulated, and several of the arguments that
depend on it --- including one in
Chapter~\ref{ch:kant} that matters a great deal --- as resting on
a construction this book has described but not exhibited.
